\section{Introduction}

% Embodied foundation models exhibit scaling laws that mirror those observed in large language models, where generalization expands commensurately with the volume of real-world interaction data~\cite{gen0,luo2025being}. Yet the supply of such data remains the primary bottleneck. Physical interaction data is expensive to acquire, difficult to diversify, and inherently noisy, making data scarcity a fundamental constraint on embodied intelligence.

% Several collection paradigms have been explored to alleviate this bottleneck, each facing a characteristic trade-off. Teleoperation~\cite{chen2025towards,droid2024} yields high-fidelity demonstrations but requires hardware exceeding \$50K and confines collection to lab settings. Universal Manipulation Interfaces~\cite{chi2024umi,xu2025dexumi} reduce per-unit cost to \$300--\$800 yet demand active, deliberate manipulation that disrupts natural behavior. Wearable AR/VR systems~\cite{wang2024dexcap,zhong2025humanoidexo} provide precise hand tracking at the cost of devices heavier than 500g, too intrusive for continuous everyday wear. Passive egocentric video~\cite{grauman2022ego4d,damen2018epic} offers scale at near-zero marginal cost but typically lacks the fine-grained manipulation annotations required for policy learning. As a result, researchers who wish to use egocentric data for embodied learning face a dual challenge: the data itself lacks manipulation-aligned geometric signals, and bridging the embodiment gap from raw video to trainable robot policies demands substantial, often idiosyncratic post-processing. Different downstream approaches, from retargeting-based transfer~\cite{qiu2025humanoid,punamiyaegobridge} to VLA pretraining~\cite{luo2025being,yang2025egovla} and world-model training~\cite{egoscale2026}, each devise their own filtering, retargeting, and format conversion pipelines, making the path from data to deployment fragmented and labor-intensive.

% Our prior work, AoE (Always-on Egocentric)~\cite{aoe2026}, introduced a collection paradigm that leverages ubiquitous smartphones for passive, always-on egocentric video acquisition at less than \$20 per user. While AoE demonstrated that consumer-grade devices can produce manipulation data of sufficient quality for policy learning, the question of how to make such data immediately usable by the broader community remained open. To this end, we present Open-AoE, an open-source release that pairs 2,000 hours of smartphone-collected egocentric manipulation data with a complete data-to-model toolchain. Open-AoE comprises four components. Open-AoE-2000H is the core dataset of 2,000 hours with English atomic action annotations, MANO hand pose reconstructions, and 6-DoF camera trajectories. AoE-Visualization provides a Rerun-based data browser for intuitive exploration and quality assessment. AoE-Retarget-Replay offers a human-to-robot motion transfer pipeline validated through real-hardware replay on Unitree G1. AoE-Training-Ready supplies format converters and training recipes for mainstream embodied models including $\pi_{0.5}$~\cite{pi05}, GR00T~\cite{groot2025}, and fastWAM~\cite{fastwam2025}.

% Open-AoE offers three distinctive advantages over existing egocentric manipulation datasets. First, the dataset achieves unprecedented sensor diversity. Data is collected across 400+ smartphone models, with calibrated horizontal FOV spanning 35.1$^{\circ}$ to 125.7$^{\circ}$. Each device-camera stack implements its own image signal processing pipeline, including white-balance, HDR, noise-reduction, color mapping, exposure, and firmware choices. This variability forces downstream models to learn camera-domain-invariant visual representations rather than exploiting device-specific shortcuts, a property that no dataset collected with a single capture device can provide. Furthermore, the collector population of 1,000+ individuals across diverse domestic and workplace environments yields rich behavioral and environmental diversity. Different operators bring distinct manipulation styles, hand shapes captured by MANO shape parameters, and uncontrolled household lighting and clutter that approximate deployment conditions far better than lab-collected alternatives.

% Second, Open-AoE provides annotation completeness through frame-aligned, multi-modal signals. Each video segment is accompanied by undistorted frames with calibrated intrinsics, MANO hand pose reconstructions comprising 45-dimensional joint rotations and 10-dimensional shape parameters per hand, 6-DoF camera trajectories with estimated depth maps, and atomic action annotations with verb, object, hand usage, bounding box, and English natural-language descriptions. All modalities are aligned by frame index, eliminating timestamp drift across annotation layers. To our knowledge, no other open-source egocentric dataset simultaneously provides all four signal types at this granularity.

% Third, Open-AoE addresses data usability by releasing a complete toolchain alongside the dataset. Most existing egocentric datasets provide data alone, leaving users to independently solve format conversion, human-to-robot retargeting, and model-specific training configuration. AoE-Retarget-Replay offers a validated retargeting pipeline that maps MANO-based hand motions to Unitree G1 joint commands, with replay demonstrations confirming that the transferred motions are physically executable. AoE-Training-Ready provides one-command format conversion to LeRobot v2.1 and pre-configured training recipes for $\pi_{0.5}$, GR00T N1.5, and fastWAM, eliminating the engineering overhead that typically dominates the first week of any new dataset adoption. By packaging both data and the toolchain to use it, Open-AoE reduces the barrier from dataset download to model training to near zero.

\begin{table}[t]
\centering
\caption{
\textbf{Comparison of egocentric manipulation datasets.} Capture devices are reported as device models or camera types only when explicitly documented by the source; aggregated corpora without a verifiable model inventory are marked Mixed / N/A. Open-AoE combines large-scale consumer-smartphone collection with complete manipulation annotations and a data-to-model toolchain.}
\label{tab:dataset_comparison}
\resizebox{\textwidth}{!}{
\begin{tabular}{lccccccccc}
\toprule
Dataset & Hours & Tasks$^{\dagger}$ & Hand Pose & Language & Camera Traj. & Participants & Capture Devices & Retarget & Training \\
\midrule
Ego4D~\cite{grauman2022ego4d} & 3,670 & N/A & \ding{55} & \ding{51} & \ding{55} & 923$^{\ast}$ & 7 cameras & \ding{55} & \ding{55} \\
EPIC-KITCHENS~\cite{damen2022rescaling} & 100 & N/A & \ding{55} & \ding{51} & \ding{55} & 37 & 2 models & \ding{55} & \ding{55} \\
EgoDex~\cite{egodex2025} & 829 & 194 & \ding{51}$^{\text{native}}$ & \ding{51} & \ding{51} & N/A & 1 model & \ding{55} & \ding{55} \\
EgoLive~\cite{egolive2026} & 1,680 & 346 & \ding{51} & \ding{51} & \ding{51} & N/A & 1 custom model & \ding{55} & \ding{55} \\
OpenEgo~\cite{openegocentric2025} & 1,107 & 290 & \ding{51}$^{\text{unified}}$ & \ding{51} & Partial & N/A & Mixed / N/A & \ding{55} & \ding{55} \\
EgoScale~\cite{egoscale2026} & 20,854 & N/A & \ding{51} & \ding{51} & \ding{55} & N/A & Mixed / N/A & \ding{55} & \ding{55} \\
\midrule
\textbf{Open-AoE} & \textbf{2,000} & \textbf{8,000+} & \ding{51}$^{\text{MANO}}$ & \ding{51} & \ding{51} & \textbf{500+} & \textbf{400+ models} & \ding{51} & \ding{51} \\
\bottomrule
\end{tabular}
}
\parbox{\textwidth}{\footnotesize
$^{\dagger}$Task counts follow source-specific definitions and are not directly comparable.
$^{\ast}$Ego4D uses the \href{https://ego4d-data.org/}{current project-page} counts (923 participants; seven cameras); its CVPR 2022 paper reports 931 camera wearers. OpenEgo's six denotes source datasets, not participants.}
\end{table}

% As shown in Table~\ref{tab:dataset_comparison}, Open-AoE is the first open-source egocentric dataset that simultaneously provides complete manipulation annotations, broad consumer-smartphone sensor coverage, and a data-to-model toolchain, collected entirely from consumer devices at less than \$20 per contributor.


% Embodied foundation models are entering a stage where capability scaling is increasingly driven by data scale. Unlike language models, which can draw from internet-scale text, robot learning depends on physical interaction data from the real world: how humans observe a scene, move their hands, contact objects, and carry out long-horizon tasks. Such data must not only be large and diverse, but also preserve the geometric, temporal, and semantic structure needed for manipulation learning. Without this kind of data, increasing model size alone is unlikely to yield robust generalization in physical tasks.

% Egocentric video offers a natural data modality for this problem. By recording interaction from the actor’s own viewpoint, it captures hands, objects, and environments in a perception structure closer to robot execution than third-person video, while being far easier to scale in everyday settings than robot teleoperation. Recent efforts such as Ego4D, EPIC-KITCHENS, EgoDex, EgoLive, and EgoScale have steadily expanded the scale and quality of egocentric data. Some provide large-scale recordings of daily activities, while others introduce hand pose, camera trajectory, or fine-grained manipulation annotations. Together, these works show that egocentric data is moving beyond visual understanding toward a data substrate for robot learning.

% Yet this line of work also reveals a more concentrated bottleneck. What the community lacks is not simply another isolated video collection, but an open data infrastructure that can continuously collect, systematically process, and directly support model training. High-quality egocentric datasets often rely on custom head-mounted cameras, AR/VR devices, or organized capture protocols. These setups provide stable viewpoints and detailed annotations, but make it difficult to build a low-cost, long-term collection ecosystem in which broad participants can contribute. In contrast, large-scale passive video is easier to obtain, but typically lacks the hand motion, camera trajectory, action segmentation, and physical annotations needed to map human behavior into robot action spaces. Moreover, even when data is publicly released, researchers must often rebuild their own pipelines for filtering, anonymization, slicing, reconstruction, retargeting, format conversion, and model integration. Releasing data does not automatically release the capability to use data.

% Representative recent works have pushed different boundaries of this problem. EgoLive demonstrates the value of high-quality real-world egocentric data for robot learning, but its path depends on custom capture hardware and centralized data production. EgoInfinity shows the potential of recovering 4D hand-object interaction from internet RGB videos and connecting it to robot actions, but focuses on retrospectively mining existing video resources rather than opening a complete production chain from real-world capture to training-ready samples. Together, these efforts point to an open question: how can low-cost real-world collection, structured post-processing, and downstream training adaptation be released as a single reusable system?

% To address this gap, we present Open-AoE, built on AoE technology. Open-AoE is not merely a dataset release; it opens the AoE data collection infrastructure to the community. Using consumer-grade smartphones as the entry point for egocentric manipulation capture, Open-AoE releases approximately 2,000 hours of real human manipulation video. It also provides the full processing stack, including on-device interaction triggering, offline quality filtering, privacy anonymization, video slicing, coarse semantic labeling, camera trajectory estimation, MANO hand reconstruction, atomic action annotation, and multi-stage quality control. Beyond data production, Open-AoE includes tools for visualization, reconstruction, cross-embodiment retargeting, and training-format conversion, allowing the same data to support VLA training, human-to-robot transfer, world action models, and world model learning. The central contribution of Open-AoE is therefore not a single annotation modality or a single model adapter, but the release of the full “capture-process-reconstruct-train” loop for real-world egocentric data.

% Concretely, the first Open-AoE release contains approximately 2,000 hours of first-person human manipulation data collected with consumer smartphones, forming a structured training corpus with video, MANO hand poses, camera trajectories, and language descriptions. For data production, we open-source the processing scripts that convert raw mobile video into anonymized training samples, covering quality filtering, video segmentation, semantic detection, reconstruction annotation, and data acceptance checks. For data use, we provide a Rerun-based visual browser, retargeting and replay pipelines for robots such as Unitree G1, and conversion scripts and training recipes for widely used frameworks including LeRobot, GR00T, π0.5, and fastWAM. Starting from the same Open-AoE corpus, researchers can move directly into data inspection, 4D reconstruction, robot replay, VLA training, or world model learning without rebuilding a private processing stack. By lowering both the cost of contributing data and the cost of using data, Open-AoE aims to turn egocentric human manipulation video from a watchable resource into trainable open infrastructure.

% Our contributions are three-fold:
% \begin{itemize}
%     \item We release approximately 2,000 hours of real egocentric human manipulation data collected with consumer-grade smartphones, providing the community with a low-cost, natural-scene, multi-device corpus for egocentric manipulation learning.
%     \item We open a complete data processing pipeline that converts raw smartphone videos into structured training data with hand poses, camera trajectories, action segments, and language descriptions, reducing the reproducibility barrier for high-quality egocentric data production.
%     \item We provide a downstream embodied learning toolchain that enables Open-AoE data to be quickly visualized, reconstructed, retargeted, and integrated into mainstream VLA, world action model, and world model training frameworks.
% \end{itemize}


Embodied foundation models are increasingly constrained by the scale and quality of real-world interaction data~\cite{lindata,droid2024}.
Unlike language models, which can learn from internet-scale text, robot learning requires physical demonstrations that capture how humans perceive scenes, move their hands, contact objects, and complete temporally extended tasks~\cite{droid2024}.
Such data must be not only large and diverse, but also structured in geometry, time, and semantics for manipulation learning~\cite{droid2024}.
Empirical scaling studies further show that diversity across environments and objects can matter more than simply increasing the number of demonstrations within fixed conditions~\cite{lindata}.

Egocentric video provides a natural and scalable modality for this problem~\cite{grauman2022ego4d,damen2022rescaling}.
By recording interaction from the actor's viewpoint, it captures hands, objects, scenes, and action progress in the same visual stream~\cite{grauman2022ego4d,damen2022rescaling}.
Compared with third-person video, egocentric video is closer to the perceptual structure of robot execution~\cite{egodex2025,egoscale2026}.
Compared with robot teleoperation, it can be collected more naturally, continuously, and cheaply in everyday environments~\cite{aoe2026,droid2024}.
Recent datasets such as Ego4D, EPIC-KITCHENS, EgoDex, EgoLive, OpenEgo, and EgoScale have advanced egocentric data along multiple axes, including larger-scale daily activity capture, finer-grained hand pose annotation, richer manipulation scenarios, and stronger alignment with robot learning~\cite{grauman2022ego4d,damen2022rescaling,egodex2025,egolive2026,openegocentric2025,egoscale2026}.

However, a central gap remains.
The community does not simply need more egocentric video.
It needs open data infrastructure that can continuously collect, systematically process, and directly support model training.
High-quality egocentric datasets often rely on specialized head-mounted devices, AR/VR hardware, or controlled capture protocols~\cite{egodex2025,egolive2026}.
Specialized hardware can improve pose and camera sensing, but it narrows accessibility relative to commodity-smartphone collection~\cite{egodex2025,egolive2026,aoe2026}.
In contrast, large-scale passive video is easier to obtain, but usually lacks the hand motion, camera trajectory, action boundaries, and physical structure required for robot training~\cite{egodex2025,openegocentric2025}.
% Recent efforts consequently introduce dataset-specific pipelines for hand-pose unification, human-to-robot transfer, or raw-video processing~\cite{openegocentric2025,egoscale2026,aoe2026}.
% Releasing videos is therefore not the same as releasing trainable data.
Recent efforts have consequently developed dataset-specific pipelines for hand-pose unification, human-to-robot transfer, and raw-video processing~\cite{openegocentric2025,egoscale2026,aoe2026}. 
However, these pipelines remain fragmented and have yet to form a unified infrastructure. This fragmentation underscores the distinction between releasing videos and releasing data that can be used directly for training.

To address this gap, we introduce \textbf{Open-AoE}, a community-oriented egocentric manipulation dataset and toolchain.
Open-AoE builds on the AoE consumer-smartphone collection framework~\cite{aoe2026} and releases the full pipeline from raw capture to model training.
The first release contains approximately \textbf{2,000 hours} of egocentric human manipulation data collected in natural environments by \textbf{500+ contributors} using \textbf{400+ smartphone models}.
The dataset provides structured signals including videos, text descriptions, MANO hand poses, camera trajectories, and atomic action segments.
For dataset production, Open-AoE releases the cloud-side processing pipeline that converts uploaded smartphone clips into structured samples through quality screening and privacy erasure, temporal segmentation, semantic labeling, camera-trajectory estimation, hand reconstruction, atomic-action annotation, and multi-stage quality inspection.
Separately, the downstream data consumption toolchain starts from the released corpus and supports synchronized visual inspection, 4D hand-object interaction reconstruction, cross-embodiment motion retargeting, and robotized-video generation.
The Training-ready toolchain then maps the aligned video, hand, camera, and language signals into action representations and training interfaces for three model directions: Vision-Language-Action (VLA) policies, World Action Models (WAMs), and World Models.

The goal of Open-AoE is not to release an isolated video collection, but to open a reusable \textbf{capture-process-reconstruct-train} data production loop.
Starting from the same egocentric corpus, researchers can inspect data, reconstruct 4D interactions, replay motions on robots, train vision-language-action models, learn world action models, or study world model learning without rebuilding a private post-processing stack.
By lowering both the barrier to contributing data and the barrier to using data, Open-AoE aims to turn low-cost smartphone-captured human manipulation video into practical open infrastructure for embodied intelligence research.

Our contributions are three-fold:
\begin{itemize}
    \item We release approximately \textbf{2,000 hours} of real egocentric human manipulation data collected with consumer-grade smartphones in natural environments, covering diverse participants, devices, and everyday manipulation scenarios.

    \item We open an end-to-end \textbf{data processing pipeline} that converts raw smartphone videos into structured training samples with hand poses, camera trajectories, action segments, and language descriptions.

    \item We provide a downstream \textbf{embodied learning toolchain} for data visualization, reconstruction, and cross-embodiment retargeting, together with training-ready representations and interfaces for three major model directions: VLA policies, WAMs, and World Models.
\end{itemize}
